\documentclass[11pt]{article}
\usepackage[T1]{fontenc}
\usepackage{lmodern,microtype}
\usepackage[a4paper,margin=25mm]{geometry}
\usepackage{amsmath,amssymb,amsthm,mathtools}
\usepackage[colorlinks=true,linkcolor=blue,urlcolor=blue,citecolor=blue]{hyperref}
\newtheorem{theorem}{Theorem}
\newtheorem{proposition}[theorem]{Proposition}
\newtheorem{corollary}[theorem]{Corollary}
\DeclareMathOperator{\Ric}{Ric}
\DeclareMathOperator{\tr}{tr}
\newcommand{\R}{\mathbb R}
\newcommand{\Id}{\mathrm{Id}}
\setlength{\emergencystretch}{2em}
\title{A complete negative Einstein strictly almost-K\"ahler four-manifold:\\verification and a connection to Fino's existence question}
\author{Research continuation note}
\date{8 October 2026}
\begin{document}
\maketitle

The explicit metric and two-form below are the construction recorded by
the user's research system, session A01, in Section~2 of
\texttt{03\_Subsequent\_Mathematical\_Developments.tex}.
This note supplies an independent full curvature calculation, the
simply connected complete extension, and an explicit connection to a
published existence question. It does not attribute the construction
to the present continuation.

\section{The geometric statement}

Let $G_0=\mathrm{SL}(2,\R)$ and let $K\cong\mathrm{SO}(2)$ be its
maximal compact subgroup, with angular parameter of period $2\pi$.
Choose left-invariant forms such that
\begin{equation}\label{eq:MC}
 d\sigma_1=-2\sigma_2\wedge\sigma_3,\qquad
 d\sigma_2=2\sigma_3\wedge\sigma_1,\qquad
 d\sigma_3=2\sigma_1\wedge\sigma_2.
\end{equation}
For $1/3<x<1$, put $D=1+3x^2$ and
\begin{align}
 G&=\frac{4D}{(1-x)^2(1+x)^2(3x-1)(3x+1)},\notag\\
 A_1&=\frac{4(3x-1)D}{3(1-x)(1+x)^2},&
 A_2&=\frac{4(3x+1)D}{3(1-x)^2(1+x)},&
 A_3&=\frac{64x^2}{(1-x)(1+x)D},\label{eq:metric-functions}\\
 X_1&=\frac{(1-x)^3(3x+1)}{16x^3},&
 X_2&=\frac{(1+x)^3(3x-1)}{16x^3}.\label{eq:X}
\end{align}
Write
\[
 g=G\,dx^2+A_1\sigma_1^2+A_2\sigma_2^2+A_3\sigma_3^2,
 \qquad e^0=\sqrt G\,dx,\quad e^i=\sqrt{A_i}\sigma_i.
\]
With $\omega_i=e^0\wedge e^i+e^j\wedge e^k$ for cyclic $(i,j,k)$, set
\begin{equation}\label{eq:Omega}
 \Omega=\sqrt{X_1}\,\omega_1+\sqrt{X_2}\,\omega_2.
\end{equation}

\begin{theorem}\label{thm:main}
The pair $(g,\Omega)$ extends smoothly over the zero section of
\[
 E=G_0\times_K\R^2\longrightarrow G_0/K\cong\mathbb H^2,
\]
where $K$ acts with weight one on $\R^2$. The resulting Riemannian
manifold is complete and diffeomorphic to $\R^4$. It satisfies
\[
 d\Omega=0,\qquad |\Omega|_g^2=2,\qquad
 \Ric_g=-3g,\qquad W_g^-=0,
\]
and
\begin{equation}\label{eq:q}
 |\nabla\Omega|^2=\frac{36(1-x^2)^3}{(1+3x^2)^3}>0
\end{equation}
at every point, including the zero section by continuous extension.
In particular this is a strictly almost-K\"ahler Einstein structure.
The metric itself is not homogeneous.
\end{theorem}

Here the orientation is the one induced by $\Omega$, so that
``self-dual'' means $W^-=0$. The scalar curvature is $-12$.
Neither neutral signature nor a conformal change of an Einstein metric
is involved.

\section{Independent curvature and closure calculation}

The calculation used here starts with the frame brackets and Koszul's
formula; it does not assume the reduced Einstein or self-duality
equations of the source manuscript. Put $dt=\sqrt G\,dx$, $v=dx/dt$,
$h_i=vH_i$, and $w_i=vV_i$. Directly from
\eqref{eq:MC}--\eqref{eq:metric-functions},
\begin{align*}
 v^2&=\frac{(1-x^2)^2(9x^2-1)}{4D},\\
 \ell:=\frac{v'}v&=-\frac{2x}{1-x^2}+
 \frac{9x}{9x^2-1}-\frac{3x}{D},\\
 H_1&=\frac12\left(\frac3{3x-1}+\frac{6x}{D}
       +\frac1{1-x}-\frac2{1+x}\right),\\
 H_2&=\frac12\left(\frac3{3x+1}+\frac{6x}{D}
       +\frac2{1-x}-\frac1{1+x}\right),\\
 H_3&=\frac1x+\frac{x}{1-x^2}-\frac{3x}{D},\\
 V_1&=-\frac{D}{4x(1+x)(3x+1)},\\
 V_2&=\frac{D}{4x(1-x)(3x-1)},\\
 V_3&=\frac{12x}{(9x^2-1)D}.
\end{align*}
For the dual orthonormal frame, write $[E_a,E_b]=v c_{ab}^{\ d}E_d$.
Its only independent nonzero coefficients are
\[
 c_{0i}^{\ i}=-H_i,\qquad c_{jk}^{\ i}=-2V_i
 \quad ((i,j,k)\text{ cyclic}).
\]
If $\nabla_{E_a}E_b=v\gamma_{ab}^{\ d}E_d$, Koszul's formula gives
\begin{equation}\label{eq:koszul}
 \gamma_{ab}^{\ d}=\tfrac12
 (c_{ab}^{\ d}-c_{bd}^{\ a}+c_{da}^{\ b}).
\end{equation}
With the convention
$R(X,Y)Z=\nabla_X\nabla_YZ-\nabla_Y\nabla_XZ-\nabla_{[X,Y]}Z$,
the full curvature follows from
\begin{align}\label{eq:full-R}
 \langle R(E_a,E_b)E_c,E_d\rangle
 =v^2\bigg[&\delta_{a0}(\gamma_{bc}^{\ d\prime}+\ell\gamma_{bc}^{\ d})
 -\delta_{b0}(\gamma_{ac}^{\ d\prime}+\ell\gamma_{ac}^{\ d})\notag\\
 &+\sum_e(\gamma_{bc}^{\ e}\gamma_{ae}^{\ d}
 -\gamma_{ac}^{\ e}\gamma_{be}^{\ d}-c_{ab}^{\ e}\gamma_{ec}^{\ d})\bigg].
\end{align}
Rational substitution yields $\Ric_{ab}=-3\delta_{ab}$.
In the orthonormal bases $(e^{0i}\pm e^{jk})/\sqrt2$, the curvature
operator has zero mixed block and
\begin{equation}\label{eq:Rblocks}
 \mathcal R^-=-\Id,
 \qquad\mathcal R^+=\operatorname{diag}(r_1,r_2,r_3),
\end{equation}
where
\begin{align*}
 r_1&=-\frac{3(x-1)^2(21x^4+18x^3+12x^2-2x-1)}{D^3},\\
 r_2&=-\frac{3(x+1)^2(21x^4-18x^3+12x^2+2x-1)}{D^3},\\
 r_3&=\frac{9(x-1)^2(x+1)^2(5x^2-1)}{D^3}.
\end{align*}
Their sum is $-3$. Thus $s=-12$ and
$W^-=\mathcal R^--(s/12)\Id=0$.

For completeness, closure can be checked without curvature:
if $f_i=\sqrt{X_i}$, the $i$th summand satisfies
\[
 d(f_i\omega_i)=
 [\dot f_i+(h_j+h_k-2w_i)f_i]e^{0jk}=0.
\]
The derivative formula
\[
 v^{-1}(\nabla_{E_a}\Omega)_{bc}
 =\delta_{a0}\Omega_{bc}'-
 \sum_e(\gamma_{ab}^{\ e}\Omega_{ec}
        +\gamma_{ac}^{\ e}\Omega_{be})
\]
then gives \eqref{eq:q} by summing over $a$ and $b<c$.
The identity $X_1+X_2=1$ gives $|\Omega|^2=2$.
A self-dual two-form of length $\sqrt2$ defines an orthogonal
almost-complex structure. Since it is closed and not parallel, that
structure is strictly almost-K\"ahler. The accompanying SymPy script
implements \eqref{eq:koszul} and \eqref{eq:full-R} directly and checks
these identities with exact rational arithmetic.

\section{Smooth extension and completeness}

Define the geodesic radial coordinate $r$ by $r=0$ at $x=1/3$ and
$dr=\sqrt G\,dx$. The exact radial equation gives analytic functions
of $r^2$ with expansions
\begin{align*}
 x-\tfrac13&=\tfrac29r^2-\tfrac2{81}r^4+O(r^6),&
 A_1&=r^2+\tfrac29r^4+O(r^6),\\
 A_2&=6+7r^2+O(r^4),&A_3&=6+7r^2+O(r^4),\\
 A_2-A_3&=\tfrac{10}{3}r^4+O(r^6),&X_2&=\tfrac83r^2+O(r^4).
\end{align*}
Analyticity follows by first using $u=\sqrt{x-1/3}$: $r$ is an odd
analytic function of $u$ with nonzero linear coefficient, and hence has
an odd analytic inverse.

Locally over $G_0/K$, take
\[
 \sigma_1=d\psi+\beta,\qquad
 \sigma_2+i\sigma_3=e^{2i\psi}\zeta_0,\qquad z=re^{i\psi}.
\]
The right $K$-action rotates the horizontal coframe with weight two.
The normal circle has period $2\pi$ and weight one. Thus
\[
 dr+ir\sigma_1=e^{-i\psi}(dz+iz\beta),\quad
 dr-ir\sigma_1=e^{i\psi}(d\bar z-i\bar z\beta).
\]
The leading normal metric is $|dz+iz\beta|^2$. Its correction has
coefficient $O(r^4)\sigma_1^2$, which is smooth because
$r^2\sigma_1=\operatorname{Im}(\bar z\,dz)+r^2\beta$ is smooth.
The horizontal trace is a smooth function of $r^2$ times
$\sigma_2^2+\sigma_3^2$. Its weight-four anisotropy has coefficient
$r^4$ times a smooth even function, and is therefore smooth as well.

To check $\Omega$, write $a_i=\sqrt{A_i}$, $a_1=rb(r^2)$ and
$x_i=\sqrt{X_i}$. Its mixed component is
\begin{align*}
 &c_+\operatorname{Re}((dr+ir\sigma_1)\wedge(\sigma_2+i\sigma_3))\\
 &\qquad+c_-\operatorname{Re}((dr-ir\sigma_1)\wedge(\sigma_2+i\sigma_3)),
 \quad c_\pm=\tfrac12x_2(a_2\pm ba_3).
\end{align*}
Here $c_+=rF(r^2)$ and $c_-=r^3H(r^2)$. The two expressions become
\[
 \operatorname{Re}[zF(r^2)(dz+iz\beta)\wedge\zeta_0],\qquad
 \operatorname{Re}[z^3H(r^2)(d\bar z-i\bar z\beta)\wedge\zeta_0],
\]
which are smooth. The remaining radial--circle and horizontal area
terms are smooth. These expressions describe the extension on the
associated disk bundle, and are compatible under changes of local
trivialization. The metric is positive definite at the bolt;
self-duality and $|\Omega|^2=2$ continue there. At the bolt,
$q=32/3>0$.

For $\Gamma=\{1\}$ the base is the contractible hyperbolic plane.
The oriented real plane bundle $E$ is therefore trivial and
$E\cong\mathbb H^2\times\R^2\cong\R^4$.
To prove completeness, let $g_B=\sigma_2^2+\sigma_3^2$ denote the
invariant complete metric on $G_0/K$. On every radial compact interval
$0\le r\le R$, the functions $A_2,A_3$ have a common positive lower
bound $c_R$. Consequently
\[
 g\ge dr^2+c_R\pi^*g_B\quad (0\le r\le R).
\]
In fact $A_2,A_3\ge6$ here, but the interval-dependent bound suffices.
As $x\uparrow1$, with $y=1-x$, one has
\[
 G\sim\frac1{2y^2},\qquad
 A_1\sim\frac8{3y},\quad A_2\sim\frac{32}{3y^2},\quad A_3\sim\frac8y.
\]
Thus $r\to\infty$. A path of finite length has bounded $r$ and
bounded base distance by the displayed metric inequality. It remains
in a disk bundle of bounded radius over a compact base ball, which is
compact. No finite-length path escapes. Equivalently, closed metric
balls are compact, and Hopf--Rinow proves completeness.

\section{Nonhomogeneity and comparison with the classical ansatz}

Since $s/12=-1$, the eigenvalues of $W^+$ are $w_i=r_i+1$.
The intrinsic metric invariant, in curvature-operator norm, is
\begin{equation}\label{eq:Wnorm}
 \tr((W^+)^2)=
 \frac{96(81x^{12}-54x^{10}+315x^8-168x^6+99x^4-18x^2+1)}{(1+3x^2)^6}.
\end{equation}
At the bolt its value is $2/3$, whereas its limit at infinity is $6$.
It is nonconstant on the regular open set as well. The metric itself,
and not merely the metric--form pair, is therefore nonhomogeneous.
In particular it is not a left-invariant metric for any transitive
Lie-group presentation.

The example also lies outside the weakly $*$-Einstein class treated
by Armstrong's ansatz~\cite{Armstrong}. In the orthonormal self-dual
basis $\omega_i/\sqrt2$, put $\sigma=\Omega/\sqrt2$. Directly,
\begin{equation}\label{eq:offdiag}
 \left|\operatorname{pr}_{\sigma^\perp}(W^+\sigma)\right|^2
 =X_1X_2(w_1-w_2)^2
 =\frac{9(1-x^2)^3(9x^2-1)}{(1+3x^2)^4}>0.
\end{equation}
Thus $\Omega$ is not a Weyl eigenform; the weakly $*$-Einstein
condition in the Einstein almost-K\"ahler case fails. This is an
explicit distinction from that classified subclass.

\section{The precise published question answered}

Fino~\cite{Fino}, immediately after Corollary~2.1 in the published
version, asks whether a real four-dimensional Lie group can admit a
non-invariant, non-Ricci-flat Einstein strictly almost-K\"ahler metric.
The question follows the nonexistence theorem for \emph{left-invariant}
Einstein strictly almost-K\"ahler structures. It contains no compactness
or completeness requirement. The published paragraph is important:
the 2003 arXiv version has the invariant nonexistence statement as
Remark~2.2, but does not contain this subsequent question.

\begin{corollary}\label{cor:Fino}
The metric--form pair in Theorem~\ref{thm:main} supplies an affirmative
answer to Fino's stated existence question. It supplies one even with
the additional requirement of completeness, on the Lie group
$(\R^4,+)$, with a non-left-invariant metric.
\end{corollary}
\begin{proof}
Choose a diffeomorphism $\R^4\to E$ and pull back the complete pair.
The Einstein and almost-K\"ahler properties are invariant under
diffeomorphism. Negative Ricci curvature excludes Ricci flatness,
\eqref{eq:q} excludes K\"ahlerness, and the nonconstant invariant
\eqref{eq:Wnorm} excludes metric homogeneity. In particular the metric
is not translation invariant on $(\R^4,+)$.
\end{proof}

There is an even shorter answer if completeness is not requested:
the regular domain $(1/3,1)\times\mathrm{SL}(2,\R)$ is diffeomorphic
to the four-dimensional Lie group $\R\times\mathrm{SL}(2,\R)$.
Pulling back the displayed formulas gives the required non-invariant
Einstein strictly almost-K\"ahler metric there. This regular-domain
metric is incomplete at the missing bolt; the complete example above
uses a different underlying manifold.

Two distinctions fix the literature comparison. First,
Apostolov--Dr\u aghici's survey~\cite{AD}, Section~6, Remark~1, discusses
the absence at that time of a negative-scalar self-dual strictly
almost-K\"ahler Einstein four-dimensional example. The present example
has exactly these properties. The \emph{specific} question in that
remark asks whether the canonical complex hyperbolic plane has such
a local almost-K\"ahler structure. That specific question was answered
negatively by Di Scala--Nagy~\cite{DSN}, Corollary~1.1: it is not the
question answered here, since \eqref{eq:Wnorm} is nonconstant and this
metric is not locally complex hyperbolic.

Second, the construction answers the stated existence question
mathematically; the present comparison does not claim it is the first
answer in the entire subsequent literature. An affirmative mathematical
example and a claim of historical priority are different assertions.

The example is noncompact, with infinite volume. Thus it does not
settle the closed negative-scalar Goldberg problem. The mathematical
conclusion established here is a complete positive-definite
four-dimensional example and its exact relationship to the recorded
non-invariant existence question.

\begin{thebibliography}{9}
\bibitem{Fino} A.~Fino,
\emph{Almost K\"ahler 4-dimensional Lie groups with $J$-invariant Ricci tensor},
Differential Geometry and its Applications \textbf{23} (2005), 26--37.
Question immediately following Corollary~2.1.
\href{https://doi.org/10.1016/j.difgeo.2005.03.003}{doi:10.1016/j.difgeo.2005.03.003}.
\bibitem{Armstrong} J.~Armstrong,
\emph{An ansatz for almost-K\"ahler, Einstein 4-manifolds},
J. Reine Angew. Math. \textbf{542} (2002), 53--84.
\href{https://doi.org/10.1515/crll.2002.009}{doi:10.1515/crll.2002.009}.
\bibitem{AD} V.~Apostolov and T.~Dr\u aghici,
\emph{The curvature and the integrability of almost-K\"ahler manifolds:
a survey}, Fields Institute Communications \textbf{35} (2003), 25--53.
\url{https://faculty.fiu.edu/~draghici/research/paperfiles/survey-web.pdf}.
\bibitem{DSN} A.~J.~Di Scala and P.-A.~Nagy,
\emph{On the uniqueness of almost-K\"ahler structures},
C. R. Math. Acad. Sci. Paris \textbf{348} (2010), 423--425.
\href{https://doi.org/10.1016/j.crma.2010.01.005}{doi:10.1016/j.crma.2010.01.005}.
\url{https://www.numdam.org/item/10.1016/j.crma.2010.01.005.pdf}.
\end{thebibliography}
\end{document}
